\phantomsection
\addcontentsline{lot}{section}{\protect\numberline{}Tabel 74 - Jumlah penderita dan kematian pada KLB}
\ra{1.3}
% Table generated by Excel2LaTeX from sheet '74'

\begin{tabular}{rY{10em}rrlllrrr}
	\multicolumn{10}{l}{Tabel 74}\\
	\multicolumn{10}{c}{JUMLAH PENDERITA DAN KEMATIAN PADA KLB MENURUT JENIS KEJADIAN LUAR BIASA (KLB)}\\
	\multicolumn{10}{c}{KABUPATEN BELITUNG TIMUR}\\
	\multicolumn{10}{c}{TAHUN \tP}\\
	\toprule
	\multirow{2}[0]{2em}{NO} & \multirow{2}[0]{=}{JENIS KEJADIAN LUAR BIASA} & \multicolumn{2}{c}{YANG TERSERANG} & \multicolumn{3}{c}{WAKTU KEJADIAN (TANGGAL)} & \multicolumn{3}{c}{JUMLAH PENDERITA}  \\
	\cmidrule(l{2pt}r{2pt}){3-4}\cmidrule(l{2pt}r{2pt}){5-7}\cmidrule(l{2pt}r{2pt}){8-10}
	& & JUMLAH KEC. & JUMLAH DESA/KEL. & DIKETAHUI & DITANGGULANGI & AKHIR & L & P & L+P \\
    \midrule
    \emph{1} & \emph{2} & \emph{3} & \emph{4} & \emph{5} & \emph{6} & \emph{7} & \emph{8} & \emph{9} & \emph{10} \\
    \midrule
    1 & Keracunan Makanan & 1 & 4 & 29-07-2025  & 29-07-2025 & 30-07-2025 & 938 & 916 & 1.854  \\
    \bottomrule
\end{tabular}%
\vspace{2ex}

\begin{tabular}{rY{10em}Z{3em}Z{3em}Z{3em}Z{3em}Z{3em}Z{3em}Z{3em}Z{3em}Z{3em}Z{3em}Z{3em}Z{3em}}
    \toprule
    \multirow{2}[0]{2em}{NO} & \multirow{2}[0]{=}{JENIS KEJADIAN LUAR BIASA} & \multicolumn{12}{c}{KELOMPOK UMUR PENDERITA} \\
    \cmidrule(l{2pt}r{2pt}){3-14}
    & & 0-7 HARI & 8-28 HARI & 1-11 BLN & 1-4 THN & 5-9 THN & 10-14 THN & 15-19 THN & 20-44 THN & 45-54 THN & 55-59 THN & 60-69 THN & 70+ THN \\
    \midrule
    \emph{1} & \emph{2} & \emph{11} & \emph{12} & \emph{13} & \emph{14} & \emph{15} & \emph{16} & \emph{17} & \emph{18} & \emph{19} & \emph{20} & \emph{21} & \emph{22} \\
    \midrule
    1 & Keracunan Makanan & 0 & 0 & 0 & 197 & 391 & 475 & 819 & 0 & 0 & 0 & 0 & 0 \\
    \bottomrule
\end{tabular}%
\vspace{2ex}

\begin{tabular}{rY{10em}rrrrrrrrrrrr}
    \toprule
    \multirow{2}[0]{2em}{NO} & \multirow{2}[0]{=}{JENIS KEJADIAN LUAR BIASA}  & \multicolumn{3}{X{6em}}{JUMLAH KEMATIAN} & \multicolumn{3}{X{11em}}{JUMLAH PENDUDUK TERANCAM} & \multicolumn{3}{X{10em}}{\emph{ATTACK RATE (\%)}} & \multicolumn{3}{X{9em}}{CFR (\%)} \\
    \cmidrule(l{2pt}r{2pt}){3-5}\cmidrule(l{2pt}r{2pt}){6-8}\cmidrule(l{2pt}r{2pt}){9-11}\cmidrule(l{2pt}r{2pt}){12-14}
    & & L & P & L+P & L & P & L+P & L & P & L+P & L & P & L+P \\
    \midrule
    \emph{1} & \emph{2} & \emph{23} & \emph{24} & \emph{25} & \emph{26} & \emph{27} & \emph{28} & \emph{29} & \emph{30} & \emph{31} & \emph{32} & \emph{33} & \emph{34} \\
    \midrule
    1 & Keracunan Makanan & 0 & 0 & 0 & 1.513 & 1.433 & 2.946 & 62,00 & 63,92 & 62,93 & 0,00 & 0,00 & 0,00 \\
    \bottomrule
\end{tabular}%

\vfill
Sumber: Subkoordinator Surveilans, Epidemiologi dan Imunisasi\par 
